\subsection{LIMITS ON THE RIGHT AND ON THE LEFT OF A FUNCTION OF A REAL VARIABLE}

Let A be a non-empty subset of \(\overline{R}\) and let \(a \neq -\infty\) be a point of \(\overline{R}\) lying in the closure of the set \(B = A \cap [-\infty, a[\) . The set B is directed with respect to the relation \(\leqslant\) , and its section filter F is the same as the trace on B of the neighbourhood filter of a in \(\overline{R}\) .

DEFINITION 2. Let \(f\) be a function defined on a non-empty subset \(A\) of \(\overline{\mathbf{R}}\) , with values in a topological space \(X\) . A limit of \(f\) with respect to the filter \(\mathfrak{F}\) , if it exists, is called a limit of \(f\) on the left at the point \(a\) , relative to \(A\) , and is denoted by

\(\lim_{x \to a, x < a, x \in \mathbf{A}} f(x), \text{ or } f(a - ), \text{ if } \mathbf{X} \text{ is Hausdorff.}\)

Likewise, if \(a\) is in the closure of the set \(A \cap ]a, +\infty]\) , we define a limit on the right (if it exists) of \(f\) at the point \(a\) , and denote it by \(\lim_{x \to a, x > a, x \in A} f(x)\) , or \(f(a+)\) , if \(X\) is Hausdorff.

The following proposition is an immediate consequence of Theorem 2:

PROPOSITION 4. Let A be a subset of \(\overline{\mathbf{R}}\) and let \(a \neq -\infty\) be a point in the closure of the intersection \(\mathbf{A} \cap [-\infty, a[\) . If \(f\) is a monotonic real-valued function defined on A, then \(f\) has a limit \(f(a - )\) on the left at \(a\) , relative to A.

